\newcommand{\ConcQoneA}{\textbf{H\textsubscript{cap}:} the MLP's train--val gap at its selected epoch is 33.6 pp, and widening it 256$\to$2048 changes training accuracy by 12.3 pp but validation accuracy by 1.9 pp. Extra capacity is spent on fitting the training set rather than on generalising, so \emph{underfitting is not the bottleneck} and H\textsubscript{cap} is rejected. \textbf{H\textsubscript{opt}:} the equal-budget deep MLP changes test accuracy by 0.7 pp. Re-shaping the same parameters into more layers does not close the gap, so H\textsubscript{opt} is rejected. \textbf{H\textsubscript{bias}:} at the same parameter count the Patch-MLP changes test accuracy by 9.3 pp (largest seed sd 0.3 pp). Under pixel permutation the Patch-MLP changes by -10.3 pp and the plain MLP by -0.2 pp, leaving the Patch-MLP -0.8 pp ahead.  All three pre-registered conditions hold: the gain comes from exploiting spatial structure, which the plain MLP cannot use. \emph{Conclusion: the missing spatial inductive bias is the main reason the MLP performs poorly, and the competing explanations are ruled out by the controls above.}}
\newcommand{\ConcQoneB}{At the validation peak (epoch 3) the corrupted training examples are predicted as their hidden true class 44.6\% of the time and as their wrong training label 8.0\% of the time (chance level for a specific wrong label given those predictions $\approx$6.2\%). Between that epoch and the end, the noisy-label fit rises by 92.0 pp to 100.0\%, true-label agreement on the noisy subset falls to 0.0\%, and validation accuracy falls by 5.4 pp, against 0.0 pp in the clean-label control. Both pre-registered conditions hold. \emph{The MLP learns generalisable class structure first (it even ``corrects'' mislabelled training images), and memorisation of the noisy labels happens afterwards, overwriting part of that structure; the drop in validation accuracy is caused by the noisy labels, since the control does not show it to the same extent.}}
\newcommand{\ConcQtwoB}{Under the pixel permutation the MLP changes by -0.2 pp and the CNN by -32.5 pp. The CNN's advantage goes from 32.5 pp on clean images to 0.2 pp. This matches the prediction. The pixels are identical, so the CNN's advantage cannot come from extra compute or easier optimisation; it disappears exactly when locality is destroyed. \emph{This supports the explanation in (a): the CNN wins because its locality and weight sharing exploit spatial structure.}}
\newcommand{\ConcQtwoC}{Accuracy drops relative to clean are 8.5 pp at 16 px, 20.7 pp at 8 px, 31.0 pp at 4 px. The drop at 16 px is only 27.4\% of the drop at 4 px: removing the global arrangement costs little as long as local content survives. \emph{E\textsubscript{local} is supported: the CNN classifies largely from local texture and colour statistics.} Errors: at 8 px the classes losing most are bird, deer, cat (27.1, 26.9, 25.8 pp) and those losing least are ship, truck; the confusion that grows most is dog$\to$cat (+91 images, seed 0). Confidence: mean confidence on wrong predictions is 80.6\% clean vs 79.0\% at 16 px, 79.8\% at 8 px, 80.9\% at 4 px; ECE is 7.4\% clean vs 12.7\%, 21.7\%, 30.3\%. The model trained on shuffled images makes its additional errors with \emph{similar or higher} confidence, i.e.\ it does not register the extra ambiguity.}
\newcommand{\ConcQtwoLink}{(a) At equal parameters the CNN leads the MLP by 32.5 pp. (b) Keeping every pixel value but destroying locality shrinks that lead to 0.2 pp, which locates the advantage in the spatial prior. (c) Shuffling blocks shows which scale of that prior matters: accuracy falls by 8.5 pp at 16 px, 20.7 pp at 8 px, 31.0 pp at 4 px. Losing the global layout (large blocks) costs much less than losing local structure (small blocks), so the CNN's advantage comes mainly from local patterns, and with 4 px blocks it is back near the MLP level. Together the three experiments locate the CNN's advantage in its spatial prior (locality and weight sharing).}
\newcommand{\ConcQthreeB}{The fixed CNN scores 88.1\% on clean images, 13.4\% on noisy images, 38.1\% on clean-AE reconstructions of noisy images and 41.2\% on DAE reconstructions. The DAE recovers 37.2\% of the clean--noisy gap. Reconstructing a \emph{clean} image through the same architecture already gives only 58.8\%, which is the ceiling any reconstruction from this model can reach. \emph{Denoising recovers the information only partially:} it removes noise (PSNR rises), but some class-relevant detail is lost, and the gap to the ceiling shows the classifier needs high-frequency detail the DAE does not restore.}
\newcommand{\ConcQthreeC}{Across the three autoencoders, the best reconstruction is AE d=1024 and the most linearly decodable class information is in AE d=1024 (probe test accuracy AE d=16: 34.1\%, AE d=128: 40.0\%, AE d=1024: 50.9\%). The DAE (53.5\%) has clean-image MSE 0.0042 vs 0.0023 for the same-size clean AE, and the random encoder control reaches 48.2\%. \emph{The two quantities agree only partly.} Among the three clean autoencoders, which differ only in capacity, lower reconstruction error does come with more class information. But the relation is not general: the DAE reconstructs clean images worse than the same-size AE yet gives the best probe accuracy, and the untrained random encoder, whose reconstruction error is far higher, still reaches 48.2\%, only 2.8 pp below the trained AE. Reconstruction rewards keeping pixel-level detail (exact colour, position, texture); class information depends on what is made linearly separable, and the denoising objective, which must ignore pixel-level noise, does that better. Much of the probe accuracy comes from the convolutional architecture itself.}
\newcommand{\ConcQfourA}{\textbf{linear}: at C=0.01 train/val are 63.3/47.0\% with 85.6\% SVs; the best C=0.01 gives 47.0\% val; at C=1000 train/val are 95.9/43.3\% with 84.8\% SVs (overfits at large C: train--val gap 52.6 pp, val 3.7 pp below best). \textbf{poly}: at C=0.01 train/val are 46.0/39.3\% with 98.3\% SVs; the best C=1 gives 51.0\% val; at C=1000 train/val are 100.0/48.6\% with 89.5\% SVs (underfits at small C: 11.7 pp below best) (overfits at large C: train--val gap 51.4 pp, val 2.4 pp below best). \textbf{rbf}: at C=0.01 train/val are 31.9/29.2\% with 100.0\% SVs; the best C=10 gives 52.3\% val; at C=1000 train/val are 100.0/51.4\% with 95.2\% SVs (underfits at small C: 23.1 pp below best) (overfits at large C: train--val gap 48.6 pp, val 0.8 pp below best). Small C tolerates many margin violations, so the boundary is too smooth and both training and validation accuracy are low (underfitting). Raising C makes violations expensive: training accuracy climbs to nearly 100\% while validation accuracy peaks and then stops improving or falls, so the extra flexibility is spent on fitting the training subset (overfitting). Unusually, the support-vector fraction stays between 80.8\% and 100.0\% for every setting and drops only slightly with C: in this 1024-d representation the classes overlap so much that almost every training point lies on or inside a margin. Fits marked $\dagger$ (linear, C$\ge$0.1) stopped at the iteration cap, so those rows are approximate; the identical numbers at the largest C values reflect this, and we do not draw conclusions from their exact values. The selected SVM (rbf, C=10) reaches 52.1\% test accuracy.}
\newcommand{\ConcQfourB}{Averaged over the 3 pairs, 92.5\% of the pair errors lie inside the margin ($|f_{ij}|<1$), and errors have smaller $|f_{ij}|$ than correct predictions (cat/dog: 0.40 vs 0.49; automobile/truck: 0.34 vs 0.57; bird/deer: 0.54 vs 0.70). The pairs need many support vectors (92.3\%, 89.5\%, 79.5\% of their training points), of which only 3.8\%, 0.2\%, 15.9\% are at the box bound C. Most support vectors are therefore \emph{free} ($0<\alpha<C$) and lie on the margin: the pair's boundary is held up by almost all of its training points rather than by a few prototypes, a sign that the two classes are nearly interleaved and the flexible RBF boundary threads between them. In representation space these pairs have low Fisher ratios (0.005, 0.014, 0.010) and nearest neighbours that frequently come from the other class (10-NN purity 54.5\%, 64.3\%, 59.8\%). The support vectors and near-boundary images in \cref{fig:q4b} share colour, background and pose across the two classes. \emph{The test errors are therefore a property of the representation, in which these classes overlap, rather than a failure of the SVM optimisation.}}
\newcommand{\ConcQfourC}{Clustering reaches test ARI 0.018 and NMI 0.050, against 52.1\% SVM accuracy on the same features. Over all 45 class pairs, the Spearman correlation between SVM confusion and cluster co-assignment is 0.37. The largest disagreement is \textbf{deer/ship}: the SVM confuses only 1.6\% of their test images, yet a random deer and a random ship image share a cluster 32.1\% of the time (within-class: 35.2\% and 32.8\%). \emph{Classification says the two classes are separable, clustering says they belong together.} The supervised margin can use a low-variance direction that separates them, whereas the unsupervised partition follows the high-variance directions (overall colour, brightness, background) along which the two classes coincide.}
